\scanpage{61}


\mspage{4}

a) $\Rightarrow$ d) trivial

d) $\Rightarrow$ c) Car \deleted{\unsure{d) si on}} $0\to N\to i_!i^*F\to F$ on aura \inserted{par d)} $0\to j^*N\to\underset{\substack{\shortparallel\\0}}{j^*i_!}i^*F$ \deleted{\unsure{exacte}} d'où $j^*N=0$, or on a \deleted{\unsure{$i^*N=0$}} aussi $i^*N=0$ [$0\to i^*N\to i^*i_!i^*F\to i^*F$] donc $N=j_*\deleted{\unsure{N|M}}H$, $H\in\mathrm{Ob}\,C''$, et comme $H=j^*\deleted{\unsure{N}}j_*H\simeq j^*N=0$, on aura $N=0$, cqfd

\margin{Cor. 2 \rule{1.2cm}{0.3pt}}

\uline{Corollaire 1} bis dual. (Les conditions \unsure{du dit|duales|ci-dessus} ne seront pas vérifiées dans les applications en vue, en particulier $j^!$ ne sera pas exact, $\operatorname{Im}i_*$ ne sera pas épaisse à droite, $F\to i_*i^*F$ ne sera pas en général surjectif].

\uline{\unsure{Remarque|Remarques}} 1 \deleted{\unsure{Si}} Les conditions du cor. 1 ne \uline{résultent pas} du fait que $i_!$ est exact, comme on voit dans la situation duale pour la localisation des modules \ldots

\uline{Remarque} 2 Sous les conditions équiv. du Cor. 1, on a la situation
\diagram{\begin{tikzcd}[ampersand replacement=\&,column sep=large]
C''\arrow[r,shift right,"j_*"] \& C\arrow[l,shift right,"j^*"']\arrow[r,shift right,"i^*"] \& C'\arrow[l,shift right,"i_!"']
\end{tikzcd}}
parfaitement \uline{symétrique} en $C',C''\deleted{\unsure{\, '|\circ}}:$
\[
\left\{\begin{aligned}
&i^*\ \text{passage au }\deleted{\text{quotient}}\ \text{exact de noyau } i_!\\
&j^*\ \text{passage au quotient exact de noyau }i_!
\end{aligned}\right.
\]
$j_*$ et $j^*$ adjoints, $i_{\deleted{\unsure{!}}}^*$ et $i_{\smash{!}}^{\deleted{\unsure{*}}}$ adjoints

De façon précise, la situation
\diagram{\begin{tikzcd}[ampersand replacement=\&,column sep=large]
C^{\prime o}\arrow[r,shift right,"i_!^{o}"] \& C^{o}\arrow[l,shift right,"i^{*o}"']\arrow[r,shift right,"j^{*o}"] \& C^{\prime\prime o}\arrow[l,shift right,"j_*^{o}"']
\end{tikzcd}}
satisfait les \deleted{\unsure{co}} \unsure{mêmes|mes|aux} hyp. que la situation de départ Cette symétrie va être détruite dans ce qui suit.

% archive sheet number: 60
